\documentclass[../main.tex]{subfiles}

\begin{document}

(Vázlat)

2.1. A „vibe coding” és az AI-asszisztált fejlesztés térnyerése
\begin{itemize}
    \item A szoftverfejlesztés paradigmaváltása: a szintaktikai precizitástól az intent-alapú fejlesztés felé.
    \item A „vibe coding” definíciója: kódírás a struktúra mély megértése nélkül, a természetes nyelvi promptok és az AI intuíciójának ötvözésével.
\end{itemize}

2.2. A kutatási probléma: A debugolás mint kritikus pont
\begin{itemize}
    \item A debugolás nem csupán hibajavítás, hanem a kód és a rendszer közötti kauzális összefüggések validálása.
    \item A probléma: Ha a kódot nem a fejlesztő alkotta meg gondolati úton, hiányzik a hibakereséshez szükséges prior knowledge.
\end{itemize}


\end{document}
